\section{Contextual equivalence and rank reduction}
\label{sec:contextual-refinement}

The labelled composition law suffices for exact recovery optimization.
This appendix asks which numerical distinctions can be removed when only
escape costs at a fixed original leaf boundary are observed.  The resulting
finite tests do not replace support families for availability or scheduling.

We first record a normal form for outer functional duals.  It will let us
replace an arbitrary outer code by the finite interface that a bounded-cost
recovery actually touches.

\begin{lemma}[Pointed column-type normal form]
\label{lem:outer-column-types}
\coverage{absent}%
Let $D=O^\perp\leq L^N$ have $L$-dimension $s$, and choose an isomorphism
$G:L^s\to D$.  There are $L$-linear coordinate covectors
$a_h\in(L^s)^\vee$ such that
\[
 G(x)_h=a_h(x)\qquad(x\in L^s, 1\leq h\leq N).
\]
The outer code projects nontrivially onto the distinguished block $j$ if and
only if the external covectors $\{a_h:h\ne j\}$ span $(L^s)^\vee$.
Conversely, every pointed list $(a_j;(a_h)_{h\ne j})$ with spanning external
part realizes such an outer context, uniquely up to change of basis in $L^s$.
\end{lemma}

\begin{proof}
Take $a_h$ to be the $h$th coordinate functional of $G$.  The target
projection of $O=D^\perp$ vanishes exactly when the target coordinate line is
contained in $D$.  Under $G$, this happens exactly when the common kernel of
the external covectors contains a vector on which $a_j$ is nonzero.  Since
$G$ is injective, the common kernel of all $a_h$ is zero, so this is equivalent
to the external covectors failing to span.  Conversely, the evaluation map
defined by a pointed list with spanning external part is injective and has no
nonzero target-supported image; taking its image as $D$ gives the required
outer code.  Changing the basis of $L^s$ changes all covectors by the same
linear automorphism and changes nothing else.
\proves{lem:outer-column-types}%
\end{proof}

For rank-one targets, the exact quotient of the labelled functions observable
by outer concatenation has a projective description.  Informally, two
rank-one inner recovery problems are indistinguishable under every compatible
outer concatenation exactly when they agree on the zero-sector cost and on the
cheapest response to every projective functional line an outer code can probe.
We now encode those probes.  Fix a line
$T=\langle u\rangle\leq W_P$.  For $b\in L$, let
$\tau_b:T\to L^*$ be defined by
\[
 \tau_b(au)(x)=a\Tr_{L/\F_q}(bx)
 \qquad(a\in\F_q,\ x\in L),
\]
and put $z_{I,T}=\rho_T(I)+d(I^\perp)$.  If
$[b]=[b_1:\cdots:b_N]\in\mathbf P(L^N)$ has a nonzero coordinate outside
the target block $j$, its projective class records an outer functional-dual
line.  Multiplying by $s\in L^\times$ changes the chosen generator of that
line but not the outer context, so the probe minimizes over this scalar choice:
\begin{equation}
 \Theta_{I,T}([b])=
 \min_{s\in L^\times}
 \left(
  \mu_{I,P,T}(\tau_{sb_j})+
  \sum_{h\ne j}\lambda_{I,T}(\tau_{sb_h})
 \right).
 \label{eq:rank-one-line-probe}
\end{equation}
Changing $u$ or the projective representative of $[b]$ does not change this
value.

Here an \emph{outer context for the rank-one target $T$} is a triple $(N,j,O)$ with $N\geq2$,
$O\leq L^N$ an $L$-linear code, and nonzero projection onto the distinguished
block $j$.  The context family ranges over every such $N$ and $j$, while the
inner target/helper split and the identified target line $T$ remain fixed.
Thus ``rank one'' refers to $\dim_{\F_q}T=1$, not to the dimension of $O$.
Operationally, two inner representations are indistinguishable at $T$ when
every compatible outer code assigns them the same exact nonconfinement cost.
The theorem below identifies exactly which numerical data determine that
response.
Two represented inner codes are \emph{contextually equivalent at $T$} when
their values of $\Gamma_{j,T}$ agree in every member of this family.  A
compatible further concatenation adds outer layers above the same
distinguished leaf.  Throughout such an evaluation, the confinement region
remains the entire original $I$-block containing the target: it is not enlarged
to a newly formed composite block.  Target coordinates and normalization also
remain fixed.  Compatibility requires that the composed outer code, viewed
over $L$, belongs to the context family just defined.

\begin{theorem}[Minimal contextual quotient at rank one]
\label{thm:rank-one-contextual-state}
\coverage{absent}%
Let $D=O^\perp\leq L^N$ under the trace identification, and suppose the
projection of $O$ onto block $j$ is nonzero.  Then
\begin{equation}
 \Gamma_{j,T}(O,I)=
 \min\left\{
 z_{I,T},
 \min_{\substack{Lb\leq D\\b\ne0}}\Theta_{I,T}([b])
 \right\},
 \label{eq:rank-one-contextual-state}
\end{equation}
where $Lb$ is the one-dimensional $L$-subspace spanned by $b$, and the inner
minimum is empty when $D=0$.

Consequently, represented inner codes over the same extension field with
identified target lines are contextually equivalent at $T$ if and only if
their values $z_{I,T}$ agree and the family of truncated profiles
\[
 [b]\longmapsto\min\{z_{I,T},\Theta_{I,T}([b])\}
\]
agrees for every $N\geq2$, target block $j$, and projective tuple with nonzero
external part.  Thus equality of this response family realizes the exact
numerical contextual quotient at rank one: no coarser numerical summary can
predict the costs in every outer context.  It is enough to test the full outer
code and outer codes whose functional dual has $L$-dimension one.  The bounded
higher-rank theorem below also shows that, after truncation at any fixed helper
radius, only finitely many of these probes are needed.
Moreover, equality of responses persists when the outer contexts are themselves
formed by compatible further concatenation, with the original leaf confinement
region fixed.
\end{theorem}

\begin{proof}
A nonzero map from the one-dimensional $\F_q$-space $T$ into $D$ is determined
by one nonzero tuple $b\in D$.  Partition $D\setminus\{0\}$ into its
$L$-lines in the nonzero-functional sector of
\eqref{eq:weighted-subspace-cost}.  Minimization on the line $Lb$ is exactly
\eqref{eq:rank-one-line-probe}, which proves
\eqref{eq:rank-one-contextual-state}.

It remains to show that every stated probe is an actual outer-code context.
The full outer code $O=L^N$ has $D=0$ and exposes $z_{I,T}$.  Given a
projective tuple $[b]$ with nonzero external part, take
$D=Lb$ and $O=D^\perp$.  Its target projection is nonzero: otherwise
$D$ would contain the target coordinate line, forcing every external
coordinate of $b$ to vanish.  This context exposes
$\min\{z_{I,T},\Theta_{I,T}([b])\}$.  Equality in every context therefore
forces equality of the displayed response family; the first part of the
theorem gives the converse.

For the closure statement, let $A$ and $C$ form a compatible outer context
$C\circ A$.  Associativity identifies
$C\circ(A\circ I_i)=(C\circ A)\circ I_i$ at the physical leaf coordinates.
On both sides, measure escape from the same original $I_i$-block, with the
same target coordinates and normalization.  Thus parenthesization preserves
the recovery equations, their helper unions, and the requirement of a nonzero
coefficient outside that region.  Since $C\circ A$ is an $L$-linear member
of the quantified context family, contextual equivalence gives equal responses
for $I_1$ and $I_2$ in this context.
\uses{thm:ungated-ranked-confinement}%
\proves{thm:rank-one-contextual-state}%
\end{proof}

This theorem compares numerical costs only.  Full coefficient witnesses still
require the minimizing lifts described after
Theorem~\ref{prop:prescribed-coset-composition}.
It does not assert congruence of escape summaries recomputed at a composite
block boundary.  In the binary code $(a,a,a,a,b,b,b,b)$ with target coordinate
$1$, escape from $\{1,2\}$ costs one helper, using $x_1+x_3=0$.
Escape from $\{1,2,3,4\}$ costs three: parity requires one helper in that
region and at least two outside it, attained by $x_1+x_2+x_5+x_6=0$.

When $\dim_{\F_q}T=t$, the image of any map
$B:T\to\operatorname{FD}(O)$ lies in an $L$-subspace
generated by at most $t$ images.  A bounded-cost map also uses only boundedly
many external labels.  Together these observations give a finite
small-context theorem.  For $c\in\mathbb N\cup\{\infty\}$ write
$[c]_{r+1}=\min\{c,r+1\}$, with $[\infty]_{r+1}=r+1$.

\begin{theorem}[Rank- and radius-bounded outer tests]
\label{prop:rank-bounded-outer-tests}
\coverage{absent}%
Let $I$ be a represented inner code with $0<\dim I<|E|$, let
$0\ne T\leq W_P$ have $\dim_{\F_q}T=t$, and let $r\geq0$.  Put
$D=O^\perp\leq L^N$, where $N\geq2$ and $O$ projects nontrivially onto block
$j$.  For $L$-subspaces $D'\leq D$, set $O_{D'}=(D')^\perp\leq L^N$.  Then
\begin{equation}
 \Gamma_{j,T}(O,I)=
 \min_{\substack{D'\leq D\\\dim_LD'\leq t}}
 \Gamma_{j,T}(O_{D'},I).
 \label{eq:rank-bounded-outer-tests}
\end{equation}
Moreover, for $S\subseteq[N]\setminus\{j\}$ put
$D_S=D\cap L^{\{j\}\cup S}$ and $O_S=D_S^\perp\leq L^N$, where $L^A$
denotes the vectors supported on $A$.  Then
\begin{equation}
 [\Gamma_{j,T}(O,I)]_{r+1}
 =\min_{\substack{S\subseteq[N]\setminus\{j\}\\|S|\leq r}}
   [\Gamma_{j,T}(O_S,I)]_{r+1}.
 \label{eq:radius-bounded-outer-tests}
\end{equation}

Consequently, if two such inner problems have different truncated responses
in some compatible outer context, a separating context exists with length at
most $\max\{2,r+1\}$ and functional-dual dimension at most
$\min\{t,r\}$.  The target leaf and target image are kept fixed under the
identifications used to compare the two problems.  The same bounds cover
witnesses: every nonconfined normalized coefficient-level system of cost at
most $r$ restricts to such a context and zero-extends back to the original
system.
\end{theorem}

\begin{proof}
For $B:T\to D$, let $D_B=\operatorname{span}_L B(T)$.  The image of an
$\F_q$-basis of $T$ generates $D_B$ over $L$, so $\dim_LD_B\leq t$.
Regarding $B$ as a map into $D_B$ changes none of its block labels or costs.
Conversely, every map into $D'\leq D$ is a map into $D$.  Minimizing the same
cost functional in both directions proves
\eqref{eq:rank-bounded-outer-tests}; $D'=0$ supplies the zero-functional
sector.  No $D'\leq D$ contains a nonzero target-supported vector, so every
$O_{D'}$ retains nonzero target projection.

For the truncated identity, a map $B:T\to D$ of cost at most $r$ has at most
$r$ nonzero external labels.  If $S$ is their set, then $B(T)\leq D_S$, and
all block costs are unchanged on passing to $O_S$.  Conversely every map into
$D_S$ is a map into $D$.  These two inclusions prove
\eqref{eq:radius-bounded-outer-tests}; maps of larger cost all truncate to
$r+1$.

Suppose now that two truncated responses differ, and choose a minimizing map
for the smaller one.  If it is in the zero-functional sector, the full outer
code of length two already separates the two inner problems.  Otherwise let
$S$ be its nonzero external-label set and let
$D_B=\operatorname{span}_L B(T)$.  Then $1\leq|S|\leq r$,
$\dim_LD_B\leq t$, and the external-coordinate projection is injective on
$D_B$, whence $\dim_LD_B\leq|S|$.  Delete every coordinate outside
$\{j\}\cup S$ and take the punctured $D_B$ as the new functional dual.
The chosen cost is unchanged.  Since this is a subcontext of the original
one, the other inner problem cannot acquire a smaller value than it had
there: every map into the punctured $D_B$ zero-extends to a map into the
original $D$.  The resulting context therefore still separates, and it has the
claimed length and functional-dual dimension.

Finally, let $Y$ be a nonconfined normalized system of cost at most $r$, put
$B_h=\Phi_IY_h$, and retain the target block together with the external blocks
on which $Y_h\ne0$.  There are at most $r$ of the latter.  Replace the
functional dual by $\operatorname{span}_L B(T)$ and delete the other blocks.
The preceding dimension argument gives the same $\min\{t,r\}$ bound.  Keeping
activity at the coefficient level, rather than retaining only blocks with
$B_h\ne0$, preserves zero-label inner-dual lifts.  Thus $Y$ restricts to the
smaller context and zero-extension recovers it exactly.
\uses{thm:ungated-ranked-confinement}%
\proves{prop:rank-bounded-outer-tests}%
\end{proof}

Over the finite fields fixed here, choose one representative of every
compatible pointed context within the two bounds in the theorem, and list its
truncated response.  This gives a finite vector
$\mathcal R_{r,T}(I)$.

\begin{corollary}[Finite coarsest bounded contextual quotient]
\label{cor:bounded-contextual-state}
\coverage{absent}%
Two represented inner target/helper problems over the same base and extension
fields, each with $0<\dim I<|E|$ and with identified target normalizations on
a nonzero space $T$, have equal values of $\mathcal R_{r,T}$ if and only if
their truncated nonconfinement costs agree in every compatible finite outer
context.
Equality of this vector is therefore the coarsest numerical equivalence that
determines exact responses through helper radius $r$.  Equality persists under
evaluation against composed outer contexts in the same family, with the
original leaf confinement region, target coordinates, and normalization fixed.

For exact, untruncated comparison, put
$z_{I,T}=\rho_T(I)+d(I^\perp)$.  This is finite for every nonzero
helper-recoverable $T$.  Two such inner problems are contextually equivalent
in all finite outer contexts if and only if their $z$ values agree, say at
$z$, and their exact responses agree on contexts of length at most
$\max\{2,z\}$ and functional-dual dimension at most
$\min\{t,z-1\}$.
\end{corollary}

\begin{proof}
The theorem says that any unequal bounded responses have a representative
among the listed contexts; the converse is immediate.  Hence equality of the
finite response vector is precisely observational equivalence, and any
numerical summary valid in all contexts must refine it.  Evaluation against a
compatible composed outer context $C\circ A$ is already among these tests.
Associativity and trace transitivity preserve the physical leaf coordinates
and target normalization; fixing the original leaf region also preserves
which equations leave it.  Equality therefore persists in these evaluations,
without asserting equality of newly formed macroblock escape summaries.

The full outer code of length two has response $z_{I,T}$ and separates unequal
$z$ values.  When both values equal $z$, every exact response is at most $z$.
Apply the theorem with $r=z-1$: equality after truncation at $z$ is then exact
equality, and the stated finite bounds follow.
\uses{lem:outer-column-types,prop:rank-bounded-outer-tests,prop:prescribed-coset-composition}%
\proves{cor:bounded-contextual-state}%
\end{proof}
